\documentclass[11pt,a4paper]{article}
\usepackage[margin=1in,landscape]{geometry}
\usepackage{booktabs}
\usepackage{pdflscape}
\usepackage{makecell}
\usepackage{placeins}
\usepackage{caption}
\usepackage{amsmath,amssymb}
\usepackage{float}
\usepackage{hyperref}
\title{\textbf{Section 3.1 IM-OLS Cointegration Estimation Matrix: U.S. NFC Transformation Elasticity (1929--2024)}}
\author{\textbf{Department of Economics, UMass Amherst}}
\date{\today}
\begin{document}
\maketitle

\section*{Data Provenance and Estimation Summary}
This report presents the empirical cointegration results for Section 3.1 of Chapter 2. All models are estimated via Cointegrating Multivariate Polynomial Regression (CMPR) IM-OLS (Vogelsang et al. 2024) over the full NIPA sample window (1931--2024, $T=94$ observations after listwise deletion for the lag-1 conditioning variables). To resolve multi-collinearity, all interactions are orthogonalized via Frisch-Waugh-Lovell (FWL) against their base regressors and any detected Bai-Perron step dummies.

\paragraph{A} \textbf{Cointegration Diagnostics and Asymptotic Null Test Behavior.} Standard residual-based tests for cointegration in non-linear polynomial regressions require careful treatment of deterministic drifts and partial sum processes. In our estimation triad, the Shin (1994) LM stationarity test ($C_{\text{Shin}}$) evaluates the null hypothesis of cointegration ($H_0$: Cointegration) directly on the Integrated Modified OLS residuals $\hat{u}_t^{+}$. Because the IM-OLS estimator is intrinsically specified in partial-sum space, the fitted residuals $\hat{u}_t^{+}$ already represent partial sums $S_t = \sum_{i=1}^t e_i$; evaluating the test without re-accumulating residuals avoids an artificial $I(2)$ divergence. Concurrently, the Vogelsang--Wagner (2014) Type 2 Wald test ($S_{\text{Type 2}}$) detrends the integrated capital regressors $x_t$ against the deterministic trend matrix $\mathbf{S}_d$ prior to constructing the reverse cumulative sum regressors $Z_{2t}$. This prevents unspanned polynomial drifts from blowing up the test statistic. Across both Specification Set 1 (double interaction) and Specification Set 2 (triple effective wage interaction), the residual ADF tests reject the null of no cointegration at standard significance thresholds once structural break step shifts ($\mathbf{D}_t^{\text{step}}$) are incorporated, while the Shin LM and VW Type 2 tests fail to reject the null of cointegration ($C_{\text{Shin}} < 0.121$, $S_{\text{Type 2}} < 13.78$), establishing robust cointegration across all preferred specifications.

\FloatBarrier
\begin{table}[H]
\centering
\caption{Model Configuration Cross-Walk}
\label{tab:model_crosswalk}
\scriptsize
\begin{tabular}{ll p{4cm} p{5.5cm} c}
\toprule
Model ID & Spec Set & Capital Measure ($k_t$) & Wage Share Measure ($\omega_{t-1}$) & Structural Breaks \\
\midrule
M1.1a & Set 1 & Private Core (ME+NRC) & Baseline NFC NVA & No \\
M1.1b & Set 1 & Private Core (ME+NRC) & Baseline NFC NVA & Yes \\
M1.2a & Set 1 & Infra-Augmented (+K\_GT) & Baseline NFC NVA & No \\
M1.2b & Set 1 & Infra-Augmented (+K\_GT) & Baseline NFC NVA & Yes \\
M1.3a & Set 1 & Private Core (ME+NRC) & Corporate NVA (Finance-Adjusted) & No \\
M1.3b & Set 1 & Private Core (ME+NRC) & Corporate NVA (Finance-Adjusted) & Yes \\
M1.4a & Set 1 & Infra-Augmented (+K\_GT) & Corporate NVA (Finance-Adjusted) & No \\
M1.4b & Set 1 & Infra-Augmented (+K\_GT) & Corporate NVA (Finance-Adjusted) & Yes \\
M1.5a & Set 1 & Private Core (ME+NRC) & Shaikh-Tonak Tier 2 (Supervisory-Corrected) & No \\
M1.5b & Set 1 & Private Core (ME+NRC) & Shaikh-Tonak Tier 2 (Supervisory-Corrected) & Yes \\
M1.6a & Set 1 & Infra-Augmented (+K\_GT) & Shaikh-Tonak Tier 2 (Supervisory-Corrected) & No \\
M1.6b & Set 1 & Infra-Augmented (+K\_GT) & Shaikh-Tonak Tier 2 (Supervisory-Corrected) & Yes \\
M2.1a & Set 2 & Private Core (ME+NRC) & Baseline NFC NVA & No \\
M2.1b & Set 2 & Private Core (ME+NRC) & Baseline NFC NVA & Yes \\
M2.2a & Set 2 & Infra-Augmented (+K\_GT) & Baseline NFC NVA & No \\
M2.2b & Set 2 & Infra-Augmented (+K\_GT) & Baseline NFC NVA & Yes \\
M2.3a & Set 2 & Private Core (ME+NRC) & Corporate NVA (Finance-Adjusted) & No \\
M2.3b & Set 2 & Private Core (ME+NRC) & Corporate NVA (Finance-Adjusted) & Yes \\
M2.4a & Set 2 & Infra-Augmented (+K\_GT) & Corporate NVA (Finance-Adjusted) & No \\
M2.4b & Set 2 & Infra-Augmented (+K\_GT) & Corporate NVA (Finance-Adjusted) & Yes \\
M2.5a & Set 2 & Private Core (ME+NRC) & Shaikh-Tonak Tier 2 (Supervisory-Corrected) & No \\
M2.5b & Set 2 & Private Core (ME+NRC) & Shaikh-Tonak Tier 2 (Supervisory-Corrected) & Yes \\
M2.6a & Set 2 & Infra-Augmented (+K\_GT) & Shaikh-Tonak Tier 2 (Supervisory-Corrected) & No \\
M2.6b & Set 2 & Infra-Augmented (+K\_GT) & Shaikh-Tonak Tier 2 (Supervisory-Corrected) & Yes \\
\bottomrule
\end{tabular}
\end{table}


\FloatBarrier
\section{Specification Set 1: Double Interactions ($k^{\text{cap}} \cdot \omega_{t-1}$)}
\FloatBarrier
\begin{table}[H]
\centering
\caption{Specification Set 1: Double Interaction FWL-Orthogonalized IM-OLS}
\label{tab:spec_set_1}
\setlength{\tabcolsep}{2.5pt}
\scriptsize
\begin{tabular}{p{3.2cm}p{1.1cm}p{1.1cm}p{1.1cm}p{1.1cm}p{1.1cm}p{1.1cm}p{1.1cm}p{1.1cm}p{1.1cm}p{1.1cm}p{1.1cm}p{1.1cm}}
\toprule
\multicolumn{13}{l}{\textbf{Dependent Variable: $y_t^{\text{NFC}}$ (Log NFC Real NVA)}} \\
\midrule
Model & M1.1a & M1.1b & M1.2a & M1.2b & M1.3a & M1.3b & M1.4a & M1.4b & M1.5a & M1.5b & M1.6a & M1.6b  \\
\midrule
$\hat{\alpha}$ (Intercept)  & 6.57$^{***}$ & 6.03$^{***}$ & 6.47$^{***}$ & 6.30$^{***}$ & 6.56$^{***}$ & 5.93$^{***}$ & 6.47$^{***}$ & 6.20$^{***}$ & 6.34$^{***}$ & 5.05$^{***}$ & 6.19$^{***}$ & 4.44$^{***}$  \\
 & {\scriptsize (0.32)} & {\scriptsize (0.73)} & {\scriptsize (0.33)} & {\scriptsize (0.68)} & {\scriptsize (0.32)} & {\scriptsize (0.76)} & {\scriptsize (0.33)} & {\scriptsize (0.71)} & {\scriptsize (0.22)} & {\scriptsize (0.75)} & {\scriptsize (0.23)} & {\scriptsize (0.68)}  \\
$\hat{\beta}$ (Private Core $k_{priv}$)  & 0.58$^{***}$ & 0.60$^{***}$ & --- & --- & 0.58$^{***}$ & 0.61$^{***}$ & --- & --- & 0.60$^{***}$ & 0.70$^{***}$ & --- & ---  \\
 & {\scriptsize (0.02)} & {\scriptsize (0.06)} &  &  & {\scriptsize (0.02)} & {\scriptsize (0.07)} &  &  & {\scriptsize (0.02)} & {\scriptsize (0.06)} &  &   \\
$\hat{\beta}$ (Infra-Augmented $k_{aug}$)  & --- & --- & 0.58$^{***}$ & 0.57$^{***}$ & --- & --- & 0.58$^{***}$ & 0.57$^{***}$ & --- & --- & 0.59$^{***}$ & 0.74$^{***}$  \\
 &  &  & {\scriptsize (0.02)} & {\scriptsize (0.06)} &  &  & {\scriptsize (0.02)} & {\scriptsize (0.06)} &  &  & {\scriptsize (0.02)} & {\scriptsize (0.05)}  \\
$\hat{\theta}$ (Orthogonalized Interaction)  & -0.02 & -0.08 & 0.02 & -0.07 & -0.03 & -0.04 & 0.01 & -0.03 & -0.24 & -0.27 & -0.24 & -0.27  \\
 & {\scriptsize (0.13)} & {\scriptsize (0.10)} & {\scriptsize (0.13)} & {\scriptsize (0.09)} & {\scriptsize (0.12)} & {\scriptsize (0.10)} & {\scriptsize (0.12)} & {\scriptsize (0.09)} & {\scriptsize (0.09)} & {\scriptsize (0.12)} & {\scriptsize (0.09)} & {\scriptsize (0.10)}  \\
Step Shifts ($D^{\text{step}}_{T_b}$)  & No & Yes & No & Yes & No & Yes & No & Yes & No & Yes & No & Yes  \\
\midrule
\multicolumn{13}{l}{\underline{Collinearity Diagnostics (VIF)}} \\
VIF ($k^{\text{cap}}$) & 1.00 & 1.00 & 1.00 & 1.00 & 1.00 & 1.00 & 1.00 & 1.00 & 1.00 & 1.00 & 1.00 & 1.00  \\
VIF (Interaction) & 1.00 & 1.00 & 1.00 & 1.00 & 1.00 & 1.00 & 1.00 & 1.00 & 1.00 & 1.00 & 1.00 & 1.00  \\
\midrule
\multicolumn{13}{l}{\underline{Cointegration Test Triad}} \\
\makecell[l]{Resid. ADF ($t_{\text{ADF}}$) \\ {\scriptsize [$H_0$: No Coint.]}} & -2.41 & -2.81 & -2.22 & -2.79 & -2.33 & -2.73 & -2.16 & -2.73 & -2.94 & -2.95 & -2.77 & -3.50$^{*}$  \\
\makecell[l]{Shin LM ($C_{\text{Shin}}$) \\ {\scriptsize [$H_0$: Coint.]}} & 0.06 & 0.04 & 0.06 & 0.04 & 0.06 & 0.04 & 0.06 & 0.04 & 0.04 & 0.03 & 0.04 & 0.03  \\
\makecell[l]{VW Type 2 ($S_{\text{Type 2}}$) \\ {\scriptsize [$H_0$: Coint.]}} & 3.30 & 11.34 & 3.74 & 10.93 & 2.70 & 11.41 & 3.22 & 10.94 & 0.05 & 11.94 & 0.11 & 10.70  \\
\midrule
Bandwidth Ratio ($b$) & 0.2746 & 0.1872 & 0.3037 & 0.1844 & 0.2864 & 0.1934 & 0.3145 & 0.1892 & 0.1940 & 0.1935 & 0.2037 & 0.1478  \\
Sample Size ($T$) & 94 & 94 & 94 & 94 & 94 & 94 & 94 & 94 & 94 & 94 & 94 & 94  \\
\bottomrule
\multicolumn{13}{p{\linewidth}}{\footnotesize Note: Estimated via Cointegrating Multivariate Polynomial Regression (CMPR) IM-OLS (Vogelsang et al. 2024) over 1929--2024. Standard errors in parentheses. Significance: $^*p<0.10, ^{**}p<0.05, ^{***}p<0.01$. Stars for coefficient rows and the Residual VW Type 2 cointegration test are derived from a custom Monte Carlo simulation suite mapping the exact DGP variant (double vs. triple interaction polynomials, with/without step level-shifts, applying FWL orthogonalization in simulation). Distributional variables centered at 1929 initial state ($t_0$).}
\end{tabular}
\end{table}


\FloatBarrier
\section{Specification Set 2: Triple Interactions ($k^{\text{cap}} \cdot \omega_{t-1} \cdot (1-\iota_{t-1})$)}
\FloatBarrier
\begin{table}[H]
\centering
\caption{Specification Set 2: Triple Interaction FWL-Orthogonalized IM-OLS}
\label{tab:spec_set_2}
\setlength{\tabcolsep}{2.5pt}
\scriptsize
\begin{tabular}{p{3.2cm}p{1.1cm}p{1.1cm}p{1.1cm}p{1.1cm}p{1.1cm}p{1.1cm}p{1.1cm}p{1.1cm}p{1.1cm}p{1.1cm}p{1.1cm}p{1.1cm}}
\toprule
\multicolumn{13}{l}{\textbf{Dependent Variable: $y_t^{\text{NFC}}$ (Log NFC Real NVA)}} \\
\midrule
Model & M2.1a & M2.1b & M2.2a & M2.2b & M2.3a & M2.3b & M2.4a & M2.4b & M2.5a & M2.5b & M2.6a & M2.6b  \\
\midrule
$\hat{\alpha}$ (Intercept)  & 6.59$^{***}$ & 6.47$^{***}$ & 6.48$^{***}$ & 6.74$^{***}$ & 6.58$^{***}$ & 6.39$^{***}$ & 6.48$^{***}$ & 6.67$^{***}$ & 6.30$^{***}$ & 5.68$^{***}$ & 6.16$^{***}$ & 5.34$^{***}$  \\
 & {\scriptsize (0.32)} & {\scriptsize (0.64)} & {\scriptsize (0.33)} & {\scriptsize (0.60)} & {\scriptsize (0.32)} & {\scriptsize (0.66)} & {\scriptsize (0.33)} & {\scriptsize (0.62)} & {\scriptsize (0.20)} & {\scriptsize (0.34)} & {\scriptsize (0.21)} & {\scriptsize (0.35)}  \\
$\hat{\beta}$ (Private Core $k_{priv}$)  & 0.58$^{***}$ & 0.56$^{***}$ & --- & --- & 0.58$^{***}$ & 0.57$^{***}$ & --- & --- & 0.60$^{***}$ & 0.65$^{***}$ & --- & ---  \\
 & {\scriptsize (0.02)} & {\scriptsize (0.06)} &  &  & {\scriptsize (0.02)} & {\scriptsize (0.06)} &  &  & {\scriptsize (0.01)} & {\scriptsize (0.03)} &  &   \\
$\hat{\beta}$ (Infra-Augmented $k_{aug}$)  & --- & --- & 0.58$^{***}$ & 0.52$^{***}$ & --- & --- & 0.58$^{***}$ & 0.53$^{***}$ & --- & --- & 0.60$^{***}$ & 0.66$^{***}$  \\
 &  &  & {\scriptsize (0.02)} & {\scriptsize (0.05)} &  &  & {\scriptsize (0.02)} & {\scriptsize (0.05)} &  &  & {\scriptsize (0.01)} & {\scriptsize (0.03)}  \\
$\hat{\theta}$ (Orthogonalized Interaction)  & 0.00 & -0.17 & 0.03 & -0.16 & -0.00 & -0.15 & 0.02 & -0.15 & -0.29$^{*}$ & -0.34$^{**}$ & -0.27$^{*}$ & -0.33$^{*}$  \\
 & {\scriptsize (0.10)} & {\scriptsize (0.07)} & {\scriptsize (0.10)} & {\scriptsize (0.07)} & {\scriptsize (0.10)} & {\scriptsize (0.07)} & {\scriptsize (0.10)} & {\scriptsize (0.07)} & {\scriptsize (0.08)} & {\scriptsize (0.08)} & {\scriptsize (0.09)} & {\scriptsize (0.08)}  \\
Step Shifts ($D^{\text{step}}_{T_b}$)  & No & Yes & No & Yes & No & Yes & No & Yes & No & Yes & No & Yes  \\
\midrule
\multicolumn{13}{l}{\underline{Collinearity Diagnostics (VIF)}} \\
VIF ($k^{\text{cap}}$) & 1.00 & 1.00 & 1.00 & 1.00 & 1.00 & 1.00 & 1.00 & 1.00 & 1.00 & 1.00 & 1.00 & 1.00  \\
VIF (Interaction) & 1.00 & 1.00 & 1.00 & 1.00 & 1.00 & 1.00 & 1.00 & 1.00 & 1.00 & 1.00 & 1.00 & 1.00  \\
\midrule
\multicolumn{13}{l}{\underline{Cointegration Test Triad}} \\
\makecell[l]{Resid. ADF ($t_{\text{ADF}}$) \\ {\scriptsize [$H_0$: No Coint.]}} & -2.32 & -3.06 & -2.17 & -3.02 & -2.26 & -2.96 & -2.12 & -2.95 & -2.96 & -3.82$^{**}$ & -2.72 & -3.78$^{**}$  \\
\makecell[l]{Shin LM ($C_{\text{Shin}}$) \\ {\scriptsize [$H_0$: Coint.]}} & 0.06 & 0.04 & 0.05 & 0.04 & 0.06 & 0.04 & 0.05 & 0.04 & 0.04 & 0.03 & 0.04 & 0.03  \\
\makecell[l]{VW Type 2 ($S_{\text{Type 2}}$) \\ {\scriptsize [$H_0$: Coint.]}} & 10.76 & 9.07 & 8.48 & 9.03 & 6.17 & 9.61 & 6.27 & 9.33 & 0.46 & 13.69 & 0.78 & 12.58  \\
\midrule
Bandwidth Ratio ($b$) & 0.2976 & 0.1687 & 0.3213 & 0.1632 & 0.3066 & 0.1732 & 0.3291 & 0.1675 & 0.1873 & 0.1651 & 0.2024 & 0.1418  \\
Sample Size ($T$) & 94 & 94 & 94 & 94 & 94 & 94 & 94 & 94 & 94 & 94 & 94 & 94  \\
\bottomrule
\multicolumn{13}{p{\linewidth}}{\footnotesize Note: Estimated via Cointegrating Multivariate Polynomial Regression (CMPR) IM-OLS (Vogelsang et al. 2024) over 1929--2024. Standard errors in parentheses. Significance: $^*p<0.10, ^{**}p<0.05, ^{***}p<0.01$. Stars for coefficient rows and the Residual VW Type 2 cointegration test are derived from a custom Monte Carlo simulation suite mapping the exact DGP variant (double vs. triple interaction polynomials, with/without step level-shifts, applying FWL orthogonalization in simulation). Distributional variables centered at 1929 initial state ($t_0$).}
\end{tabular}
\end{table}


\end{document}
